% (User input window)
% 2026-02-15 05:21:32（Asia/Singapore）

\begin{conclusion}[\texorpdfstring{$(q=2)$}{(q=2)} 的重量—碰撞配分函数：六阶常系数递推与有理生成函数]\label{con:xi-terminal-zm2-order6-rational}
令
\(
X_m=\{x\in\{0,1\}^m:\ x_ix_{i+1}=0\}
\)，
并记 Fold 纤维多重度 \(d_m(x)=|\Fold_m^{-1}(x)|\)。
定义二阶重量—碰撞配分函数
\[
Z_m^{(2)}(y):=\sum_{x\in X_m} y^{|x|_1}\,d_m(x)^2\in\ZZ[y].
\]
则对一切 \(m\ge 7\) 成立严格六阶递推
\[
\begin{aligned}
Z_m^{(2)}(y)=&\ Z_{m-1}^{(2)}(y)+(2+3y)Z_{m-2}^{(2)}(y)+(y-2)Z_{m-3}^{(2)}(y)\\
&\ +(-1-2y-3y^2)Z_{m-4}^{(2)}(y)+(1-y)Z_{m-5}^{(2)}(y)+(y+y^3)Z_{m-6}^{(2)}(y),
\end{aligned}
\]
并且其常规生成函数为显式有理函数
\[
\boxed{\;
\sum_{m\ge 0} Z_m^{(2)}(y)t^m
=\frac{1+yt+(1-2y)t^2+(y-2y^2)t^3+(y^2-y)t^4+(y^3-y^2)t^5}{1-t-(2+3y)t^2+(2-y)t^3+(1+2y+3y^2)t^4-(1-y)t^5-(y+y^3)t^6}\; }.
\]
其中可取初值
\[
Z_0^{(2)}(y)=1,\ 
Z_1^{(2)}(y)=y+1,\ 
Z_2^{(2)}(y)=2y+4,\ 
Z_3^{(2)}(y)=y^2+9y+4,
\]
\[
Z_4^{(2)}(y)=6y^2+21y+9,\ 
Z_5^{(2)}(y)=y^3+31y^2+47y+9,\ 
Z_6^{(2)}(y)=15y^3+96y^2+93y+16.
\]
\end{conclusion}
\begin{proof}[证明要点]
与结论 \ref{con:xi-terminal-zm-order4-rational} 同理，
\(Z_m^{(2)}(y)\) 是有限状态规约器的二阶碰撞计数在重量变量 \(y^{|x|_1}\) 下的扭曲求和，
故其 \(t\)-生成函数为有理函数，约分后分母次数等于最小线性实现的维数，从而诱导常系数递推。
分子分母的显式系数可由有限窗口精确枚举与有理重建（Pad\'e/Prony）恢复，并可由脚本
\texttt{scripts/exp\_fold\_zq\_weight\_collision\_partition\_audit.py} 一键复核。证毕。
\end{proof}

\begin{conclusion}[\texorpdfstring{$(q=3)$}{(q=3)} 的重量—碰撞配分函数：八阶常系数递推与有理生成函数]\label{con:xi-terminal-zm3-order8-rational}
定义三阶重量—碰撞配分函数
\[
Z_m^{(3)}(y):=\sum_{x\in X_m} y^{|x|_1}\,d_m(x)^3\in\ZZ[y].
\]
则对一切 \(m\ge 9\) 成立严格八阶递推
\[
\begin{aligned}
Z_m^{(3)}(y)=&\ Z_{m-1}^{(3)}(y)+(3+4y)Z_{m-2}^{(3)}(y)+(4y-3)Z_{m-3}^{(3)}(y)\\
&\ +(-3+y-6y^2)Z_{m-4}^{(3)}(y)+(3-6y+y^2)Z_{m-5}^{(3)}(y)\\
&\ +(1+2y-5y^2+4y^3)Z_{m-6}^{(3)}(y)+(-1+2y-y^2)Z_{m-7}^{(3)}(y)\\
&\ +(-y^4+y^3+y^2-y)Z_{m-8}^{(3)}(y),
\end{aligned}
\]
并且生成函数可写为
\[
\sum_{m\ge 0} Z_m^{(3)}(y)t^m=\frac{N_3(t,y)}{D_3(t,y)},
\]
其中
\[
N_3(t,y)=1+yt+(4-3y)t^2+(4y-3y^2)t^3+(1-6y+3y^2)t^4+(y-6y^2+3y^3)t^5-y(y-1)^2t^6-y^2(y-1)^2t^7,
\]
并且
\[
\begin{aligned}
D_3(t,y)=&\ 1-t-(3+4y)t^2-(4y-3)t^3+(6y^2-y+3)t^4-(y^2-6y+3)t^5\\
&\ -(4y^3-5y^2+2y+1)t^6+(y-1)^2t^7+(y^4-y^3-y^2+y)t^8.
\end{aligned}
\]
其中可取初值
\[
Z_0^{(3)}(y)=1,\ 
Z_1^{(3)}(y)=y+1,\ 
Z_2^{(3)}(y)=2y+8,\ 
Z_3^{(3)}(y)=y^2+17y+8,
\]
\[
Z_4^{(3)}(y)=10y^2+51y+27,\ 
Z_5^{(3)}(y)=y^3+79y^2+153y+27,
\]
\[
Z_6^{(3)}(y)=37y^3+324y^2+395y+64,\ 
Z_7^{(3)}(y)=y^4+298y^3+1304y^2+837y+64,\ 
Z_8^{(3)}(y)=101y^4+1755y^3+3965y^2+1822y+125.
\]
\end{conclusion}
\begin{proof}[证明要点]
同上：有限状态碰撞计数在重量变量下保持有限维线性实现，
从而给出有理生成函数与常系数递推；显式系数与初值由有限数据重建与枚举核验得到，
并由脚本 \texttt{scripts/exp\_fold\_zq\_weight\_collision\_partition\_audit.py} 输出可审计证书。证毕。
\end{proof}

\begin{conclusion}[\texorpdfstring{$(q=4)$}{(q=4)} 的重量—碰撞配分函数：十阶常系数递推与核多项式的在位回收]\label{con:xi-terminal-zm4-order10-rational}
定义四阶重量—碰撞配分函数
\[
Z_m^{(4)}(y):=\sum_{x\in X_m} y^{|x|_1}\,d_m(x)^4\in\ZZ[y].
\]
则对一切 \(m\ge 11\) 成立严格十阶递推
\[
\begin{aligned}
Z_m^{(4)}(y)=&\ Z_{m-1}^{(4)}(y)+(4+5y)Z_{m-2}^{(4)}(y)+(11y-4)Z_{m-3}^{(4)}(y)\\
&\ +(-10y^2+19y-6)Z_{m-4}^{(4)}(y)+(11y^2-18y+6)Z_{m-5}^{(4)}(y)\\
&\ +(10y^3-27y^2+2y+4)Z_{m-6}^{(4)}(y)+(y^3-13y^2+9y-4)Z_{m-7}^{(4)}(y)\\
&\ +(-5y^4+5y^3+2y^2-3y-1)Z_{m-8}^{(4)}(y)+(-y^3+2y^2-2y+1)Z_{m-9}^{(4)}(y)\\
&\ +(y^5-y^4+2y^3-y^2+y)Z_{m-10}^{(4)}(y),
\end{aligned}
\]
并且生成函数 \(F_4(t,y)=\sum_{m\ge 0}Z_m^{(4)}(y)t^m\) 为有理函数 \(N_4(t,y)/D_4(t,y)\)，其中分母由递推系数给出
\[
D_4(t,y)=1-\sum_{j=1}^{10}a_j(y)t^j
\]
（\(a_j(y)\) 为上式右端的十个系数多项式），而分子可取显式整系数形式
\[
\begin{aligned}
N_4(t,y)=&\ 1+yt+(11-4y)t^2+(11y-4y^2)t^3+(11-18y+6y^2)t^4\\
&\ +(11y-18y^2+6y^3)t^5+(1-13y+9y^2-4y^3)t^6+(y-13y^2+9y^3-4y^4)t^7\\
&\ +y(y-1)(y^2-y+1)t^8+y^2(y-1)(y^2-y+1)t^9.
\end{aligned}
\]
特别地，取 \(y=1\)（忘却重量层，仅保留四阶碰撞矩）时，分母在 \(t\)-变量上精确析出不可约核因子
\[
2t^5-2t^4-7t^2-2t+1,
\]
并且该因子在 \(F_4(t,1)\) 中保持在位（不被分子约去），从而以代数不可约方式回收为四阶碰撞核的本征结构常数。
\end{conclusion}
\begin{proof}[证明要点]
有理性与递推由有限状态碰撞核的有限维线性实现给出；
分子分母的显式形式及 \(y=1\) 的核因子析出可由脚本
\texttt{scripts/exp\_fold\_zq\_weight\_collision\_partition\_audit.py} 在 \(\ZZ[y,t]\) 中精确化简并输出证书。证毕。
\end{proof}

\begin{remark}[\texorpdfstring{$y=1$}{y=1} 的幽灵模态全消与 moment-kernel 的严格同态化]\label{rem:xi-terminal-zq-ghost-cancellation-y1}
记 \(F_q(t,y)=\sum_{m\ge 0}Z_m^{(q)}(y)t^m=N_q(t,y)/D_q(t,y)\)。
在 \(y=1\) 处出现统一的“幽灵模态”严格约去：
\begin{enumerate}
  \item \(q=2\)：\(N_2(t,1)=-(t-1)(t+1)^2\)，且 \(D_2(t,1)=-(t-1)(t+1)^2(2t^3-2t^2-2t+1)\)，故
  \[
  F_2(t,1)=\frac{1}{2t^3-2t^2-2t+1}.
  \]
  \item \(q=3\)：\(N_3(t,1)=-(t-1)(t+1)^2(2t^2+1)\)，且 \(D_3(t,1)=-(t-1)(t+1)^2(2t^3-4t^2-2t+1)\)，故
  \[
  F_3(t,1)=\frac{2t^2+1}{2t^3-4t^2-2t+1}.
  \]
  \item \(q=4\)：\(N_4(t,1)=-(t-1)(t+1)^2(t^2+1)(7t^2+1)\)，且
  \(D_4(t,1)=-(t-1)(t+1)^2(t^2+1)(2t^5-2t^4-7t^2-2t+1)\)，故
  \[
  F_4(t,1)=\frac{7t^2+1}{2t^5-2t^4-7t^2-2t+1}.
  \]
\end{enumerate}
因此在 \(y=1\) 的“碰撞投影”下，所有额外的 \(\pm1\) 与 \(\pm i\) 幽灵模态均被分子严格抵消，
留下的唯一动力学分母恰为 moment-kernel 的核多项式；该约去机制把重量细分与碰撞核谱系在生成函数层面严格同态化。
\end{remark}

\begin{conclusion}[谱分歧判别式的结构分解：普适因子与割圆因子的出现]\label{con:xi-terminal-zq-discriminant-factorization}
令 \(\Pi_q(\lambda,y)\) 表示 \(q=2,3,4\) 的递推特征多项式（即 \(D_q(t,y)\) 在 \(t=1/\lambda\) 下清分母后得到的 \(\lambda\)-多项式）。
则其关于 \(\lambda\) 的判别式在 \(\ZZ[y]\) 中具有严格分解：
\begin{enumerate}
  \item \(q=2\)：
  \[
  \mathrm{Disc}_{\lambda}(\Pi_2)=4\,y^{3}(y-1)\,P_{9}(y),
  \]
  其中
  \[
  \begin{aligned}
  P_{9}(y)=&\ 186624y^9-695632y^8+980411y^7-1121640y^6+1021286y^5\\
  &-458803y^4+127602y^3-35316y^2-10260y+6912.
  \end{aligned}
  \]
  \item \(q=3\)：
  \[
  \mathrm{Disc}_{\lambda}(\Pi_3)=-16\,y^{5}(y-1)^3\,P_{17}(y),
  \]
  其中 \(P_{17}(y)\in\ZZ[y]\) 为次数 \(17\) 多项式，且具有唯一负实根 \(y\approx-1.07316827\)。
  \item \(q=4\)：
  \[
  \mathrm{Disc}_{\lambda}(\Pi_4)=9216\,y^{7}(y-1)(y^2-y+1)\,R_{31}(y),
  \]
  其中 \(R_{31}(y)\in\ZZ[y]\) 为次数 \(31\) 多项式。
\end{enumerate}
因此除去普适因子 \(y^{2q-1}\) 与 \(y=1\) 的投影退化外，
\(q=4\) 首次强制出现割圆因子 \(y^2-y+1\)（即 \(6\) 次单位根），
表明四阶重量细分在谱分歧层面已不可约地感知到单位圆对称通道的额外算术结构。
\end{conclusion}
\begin{proof}[证明要点]
将 \(\Pi_q(\lambda,y)\) 视为 \(\lambda\)-多项式并计算判别式，在 \(\ZZ[y]\) 中作因式分解即可得到分解结构；
剩余因子 \(P_{17},R_{31}\) 的显式系数可由同一分解过程唯一确定，并由脚本
\texttt{scripts/exp\_fold\_zq\_weight\_collision\_partition\_audit.py} 导出为可审计工件。证毕。
\end{proof}

\begin{conclusion}[零重量纤维的完全闭式：\texorpdfstring{$d_m(0^m)=\lfloor (m+2)/2\rfloor$}{d_m(0^m)=floor((m+2)/2)}]\label{con:xi-terminal-zero-weight-fiber-closed}
记 \(0^m\in X_m\) 为全零稳定类型。
则其纤维大小满足对一切 \(m\ge 1\) 的精确公式
\[
d_m(0^m)=\Bigl\lfloor\frac{m+2}{2}\Bigr\rfloor.
\]
从而对任意整数 \(q\ge 1\) 有
\[
Z_m^{(q)}(0)=d_m(0^m)^q=\Bigl\lfloor\frac{m+2}{2}\Bigr\rfloor^q,
\]
给出重量细分在 \(k=0\) 层上的完全可解子扇区；该层的增长为多项式阶而非指数阶，对应一个严格的“边界零熵”极端相。
\end{conclusion}
\begin{proof}
由 Fold 的定义（见脚本 \texttt{scripts/common\_phi\_fold.py} 的 Zeckendorf 正规形口径），
\(\Fold_m(a)=0^m\) 当且仅当微态数值 \(N(a)\) 的 Zeckendorf 数码在位置 \(1,\dots,m\) 全为 \(0\)；
等价地 \(N(a)\in\{0,F_{m+2}\}\)。
其中 \(N(a)=0\) 仅对应全零微态，贡献 \(1\)。
另一方面，设 \(r_m\) 为用权集合 \(\{F_2,\dots,F_{m+1}\}\) 表示 \(F_{m+2}\) 的 \(0/1\) 子集表示数。
由于 \(\sum_{i=2}^{m}F_i=F_{m+2}-1\)，任意表示必须包含 \(F_{m+1}\)，故 \(r_m\) 等于用 \(\{F_2,\dots,F_m\}\) 表示 \(F_m\) 的表示数。
对后者再作“最大项是否取用”的分解：若取用 \(F_m\) 得到一项表示；否则由于 \(\sum_{i=2}^{m-2}F_i=F_m-1\)，表示必包含 \(F_{m-1}\)，余项为 \(F_{m-2}\)，从而表示数满足递推 \(r_m=1+r_{m-2}\)（\(m\ge 3\)），且 \(r_1=0,r_2=1\)。
解得 \(r_m=\lfloor m/2\rfloor\)，因此
\(
d_m(0^m)=1+r_m=\lfloor (m+2)/2\rfloor
\)。
代回 \(Z_m^{(q)}(0)\) 即得结论。证毕。
\end{proof}

\begin{conclusion}[\texorpdfstring{$t=\pm1$}{t=±1} 直线的零轨迹：重量参数的代数曲线量子化]\label{con:xi-terminal-zq-line-tpm1-zeros}
设 \(D_q(t,y)\) 为 \(F_q(t,y)=N_q(t,y)/D_q(t,y)\) 的分母多项式。
则在 \(q=2,3,4\) 上均有严格等式
\[
D_q(1,y)=D_q(-1,y),
\]
并且其零点集合在 \(y\)-平面上量子化为有限代数集合：
\[
D_2(\pm1,y)=-y(y-1)(y-2),
\]
\[
D_3(\pm1,y)=y(y-1)(y^2-4y+6),
\]
\[
D_4(\pm1,y)=-y(y-1)(y^3-5y^2+12y-24).
\]
因此，“出现特征模态 \(\lambda=\pm1\)”仅在 \(y\) 落入上述代数集合时发生。
并且在上述三式中，除去共同因子 \(y(y-1)\) 后出现的附加因子在常数项处满足
\[
(y-2)\big|_{y=0}=-2,\qquad
(y^2-4y+6)\big|_{y=0}=6,\qquad
(y^3-5y^2+12y-24)\big|_{y=0}=-24,
\]
即呈现统一的阶乘尺度律 \(({-}1)^{q+1}q!\)（对 \(q=2,3,4\)）。
\end{conclusion}
\begin{proof}[证明要点]
把 \(t=\pm1\) 代入各 \(D_q(t,y)\) 并在 \(\ZZ[y]\) 中因式分解即可；
等式 \(D_q(1,y)=D_q(-1,y)\) 反映分母在 \(t\mapsto -t\) 的偶性结构在该三阶族中的稳定性。证毕。
\end{proof}

